% AUTO-GENERATED by scripts/exp_real_input_40_output_potential_ldp_algebraic_param.py
\paragraph{输出位势的大偏差：代数参数化（可复现）}
令 $u=e^{\theta}$ 且 $\lambda(u)=\rho(M(u))$ 为 $G(\lambda,u)=0$ 的最大正实根（命题 \ref{prop:real-input-40-pressure-freezing}）。
则平衡态下输出 $1$ 的典型密度满足
$$
\alpha(u):=\frac{d}{d\theta}\log\lambda(e^{\theta})=\frac{u\,\lambda'(u)}{\lambda(u)}=-\frac{u\,\partial_u G(\lambda,u)}{\lambda\,\partial_{\lambda}G(\lambda,u)}\Bigg|_{\lambda=\lambda(u)}.
$$
对本核的显式 $G$ 展开偏导后，上式可写为完全有理形式（其中分子分母均为多项式）：
$$
\begin{aligned}
\alpha(u)&=\left.\frac{N(\lambda,u)}{D(\lambda,u)}\right|_{\lambda=\lambda(u)}.\\
N(\lambda,u)&=- 4 u^{4} + 3 u^{3} + \lambda^{3} u - \lambda^{4} u - 6 \lambda u^{2}\\
&\qquad - 6 \lambda^{6} u - 6 \lambda^{3} u^{2} - 4 \lambda^{2} u^{2}\\
&\qquad - 3 \lambda^{2} u^{3} + 2 \lambda^{5} u + 12 \lambda u^{3}\\
&\qquad + 18 \lambda^{4} u^{2}\\
D(\lambda,u)&=- 8 \lambda^{8} + 7 \lambda^{7} - 36 \lambda^{4} u^{2} - 10 \lambda^{5} u\\
&\qquad - 4 \lambda u^{3} - 3 \lambda^{3} u + 2 \lambda^{2} u^{3} + 3 \lambda u^{2}\\
&\qquad + 4 \lambda^{4} u + 4 \lambda^{2} u^{2} + 9 \lambda^{3} u^{2}\\
&\qquad + 36 \lambda^{6} u
\end{aligned}
$$
相应的大偏差率函数（Legendre 对偶）沿同一参数 $u$ 具有显式参数表示：
$$
I(\alpha(u))=\alpha(u)\log u-\log\frac{\lambda(u)}{\lambda(1)}.
$$
由于同时满足 $G(\lambda,u)=0$ 与上式的有理约束，故 $(\alpha,u)$ 或 $(\alpha,\lambda)$ 之间可通过消元得到一条代数曲线，从而端点展开与代数误差界均可在符号层面审计。
